% Options for packages loaded elsewhere
\PassOptionsToPackage{unicode}{hyperref}
\PassOptionsToPackage{hyphens}{url}
%
\documentclass[
  10pt,
]{article}
\usepackage{amsmath,amssymb}
\usepackage{iftex}
\ifPDFTeX
  \usepackage[T1]{fontenc}
  \usepackage[utf8]{inputenc}
  \usepackage{textcomp} % provide euro and other symbols
\else % if luatex or xetex
  \usepackage{unicode-math} % this also loads fontspec
  \defaultfontfeatures{Scale=MatchLowercase}
  \defaultfontfeatures[\rmfamily]{Ligatures=TeX,Scale=1}
\fi
\usepackage{lmodern}
\ifPDFTeX\else
  % xetex/luatex font selection
\fi
% Use upquote if available, for straight quotes in verbatim environments
\IfFileExists{upquote.sty}{\usepackage{upquote}}{}
\IfFileExists{microtype.sty}{% use microtype if available
  \usepackage[]{microtype}
  \UseMicrotypeSet[protrusion]{basicmath} % disable protrusion for tt fonts
}{}
\makeatletter
\@ifundefined{KOMAClassName}{% if non-KOMA class
  \IfFileExists{parskip.sty}{%
    \usepackage{parskip}
  }{% else
    \setlength{\parindent}{0pt}
    \setlength{\parskip}{6pt plus 2pt minus 1pt}}
}{% if KOMA class
  \KOMAoptions{parskip=half}}
\makeatother
\usepackage{xcolor}
\usepackage[margin=2cm]{geometry}
\usepackage{longtable,booktabs,array}
\usepackage{calc} % for calculating minipage widths
% Correct order of tables after \paragraph or \subparagraph
\usepackage{etoolbox}
\makeatletter
\patchcmd\longtable{\par}{\if@noskipsec\mbox{}\fi\par}{}{}
\makeatother
% Allow footnotes in longtable head/foot
\IfFileExists{footnotehyper.sty}{\usepackage{footnotehyper}}{\usepackage{footnote}}
\makesavenoteenv{longtable}
\setlength{\emergencystretch}{3em} % prevent overfull lines
\providecommand{\tightlist}{%
  \setlength{\itemsep}{0pt}\setlength{\parskip}{0pt}}
\setcounter{secnumdepth}{-\maxdimen} % remove section numbering
\ifLuaTeX
  \usepackage{selnolig}  % disable illegal ligatures
\fi
\IfFileExists{bookmark.sty}{\usepackage{bookmark}}{\usepackage{hyperref}}
\IfFileExists{xurl.sty}{\usepackage{xurl}}{} % add URL line breaks if available
\urlstyle{same}
\hypersetup{
  hidelinks,
  pdfcreator={LaTeX via pandoc}}

\author{}
\date{}

\begin{document}

\hypertarget{can-meaning-have-universal-coordinates-a-negative-and-positive-result}{%
\section{Can Meaning Have Universal Coordinates? A Negative-and-Positive
Result}\label{can-meaning-have-universal-coordinates-a-negative-and-positive-result}}

\emph{Draft for arXiv (cs.CL; cross-list cs.IR). The Sphota Authors,
2026. Not yet submitted.}

\hypertarget{abstract}{%
\subsection{Abstract}\label{abstract}}

We tried to build a \emph{universal coordinate system for meaning}: a
frozen, open, hierarchical code (e.g.~\texttt{JYZA-7M-X0-A8}) that any
text in any language would receive, reachable from many embedding models
through linear adapters, so that translations, paraphrases and different
models agree on the same address. Before running the decisive
experiments we fixed a four-part pass/fail gate on held-out data
(coverage, translation agreement, cross-model agreement, language
collapse), and later amended it to close a loophole: the gate must be
met with at least 256 top-level cells. \textbf{No configuration we built
passes.} The frozen release (4096 top-level cells, multilingual-e5-small
anchor) fails all four criteria (coverage 29\%, confident translation
agreement 65\%, cross-model 16\%, worst language collapse 34\%).
Agreement is bought with coarseness: each 4× finer top level costs
roughly 10--25 points on the gate metrics, and every near-pass uses
32--64 cells. The anchor matters (LaBSE beats e5-small by 11 points of
coverage and 4.5 points of confident translation agreement at the same
granularity) but not enough to change the conclusion: a GPU run at
\textasciitilde450k sentences with LaBSE and bge-m3 anchors removes
language collapse and reaches 88\% confident translation agreement at
256 cells, yet cross-model agreement peaks at 50\% (bar: 60\%). We also
document a \emph{language-collapse} failure: at ≥ 256 cells, k-means
forms cells that are mostly one language (up to 55\% of Portuguese test
sentences), which a diversified sample only partly fixes. What does
work: (1) \emph{coarse regions} (32--64 cells) where translations agree
78--89\% of the time when confident, useful for multilingual topic
analytics, routing and caching; and (2) \emph{migration through stored
canonical vectors}: an index built with one model can be queried with
another at 96.5--99.1\% recall@10 without re-embedding. We release
everything, including the gate, the negative results and the scripts to
rerun them.

\hypertarget{introduction}{%
\subsection{1. Introduction}\label{introduction}}

Embedding vectors are a good \emph{computational} representation of
meaning and a poor \emph{data format}: a vector from model A means
nothing to model B, it cannot be read or grouped in SQL, and every model
upgrade forces re-embedding the corpus. A natural ambition is a shared,
stable \emph{address} for meaning, like latitude and longitude:
readable, hierarchical, frozen, and the same whichever language or model
produced the text.

This paper asks whether that ambition is achievable with current
multilingual encoders and simple, transparent machinery (linear
projections, residual k-means, linear adapters). Our answer is split:

\begin{itemize}
\tightlist
\item
  \textbf{Negative.} A \emph{fine-grained} universal code (hundreds to
  thousands of top-level cells, stable across languages and models) is
  not supported by our measurements. The failure is structural, not a
  tuning issue: any partition of a continuous space splits equivalent
  texts near its boundaries, and the fraction split grows quickly with
  the number of cells.
\item
  \textbf{Positive.} A \emph{coarse} shared partition is stable enough
  to be useful, and the continuous canonical space behind it is an
  effective hub for cross-model migration.
\end{itemize}

We think the negative result is worth publishing on its own. ``Semantic
IDs'' and shared embedding spaces are active topics, and claims of a
universal meaning code are easy to make on in-sample numbers; ours
looked good in-sample too (81\% of training sentences get a confident
top-level code versus 26\% held out).

The project name comes from Bhartṛhari's \emph{sphoṭa} (5th c.~CE):
meaning that bursts forth whole and stays the same across its
utterances. Our results suggest it stays the same only at a coarse
resolution.

\textbf{Contributions.} 1. A complete, open method for a frozen
hierarchical meaning code (§3). 2. A pre-registered decision gate, plus
an amendment closing a granularity loophole (§4). 3. Full gate results
across granularity, sample design and anchor (§5). 4. The granularity
trade-off, the anchor effect and the language-collapse failure (§5--6).
5. Evidence for what does work: coarse regions and vector-based
migration (§7).

\hypertarget{related-work}{%
\subsection{2. Related work}\label{related-work}}

\textbf{Semantic IDs.} TIGER (Rajput et al., 2023) quantizes item
embeddings with an RQ-VAE into codeword tuples for generative retrieval;
DSI (Tay et al., 2022) uses hierarchical k-means document IDs; Singh et
al.~(2024) study Semantic IDs for ranking. These codes are private to
one system and refit per catalogue, so they never face the
cross-language and cross-model stability test that a \emph{universal}
code must pass. Codebook collapse in RQ-VAE is mitigated by EMA updates
and dead-code revival (Zeghidour et al., 2021; Lee et al., 2022).

\textbf{Embedding translation.} vec2vec (Jha et al., 2025) translates
between embedding spaces without paired data, in support of the Platonic
representation hypothesis (Huh et al., 2024). Our results are compatible
with that hypothesis at the level of \emph{continuous} geometry
(migration works) while showing that \emph{discrete} agreement is far
weaker.

\textbf{Language-agnostic spaces.} LASER, LaBSE (Feng et al., 2022) and
SONAR (Duquenne et al., 2023) embed many languages jointly; Large
Concept Models (2024) operate in SONAR space. We quantify how much of
their cross-lingual alignment survives discretization.

\textbf{Topic models.} BERTopic (Grootendorst, 2022) clusters embeddings
and labels clusters with c-TF-IDF, refitting per corpus. Our positive
result is close to a \emph{frozen, multilingual, shared} BERTopic-style
partition.

\hypertarget{method}{%
\subsection{3. Method}\label{method}}

\hypertarget{anchor-encoder}{%
\subsubsection{3.1 Anchor encoder}\label{anchor-encoder}}

\texttt{intfloat/multilingual-e5-small} (MIT, 384-d, \textasciitilde100
languages), input \texttt{"query:\ "\ +\ text}, mean pooling, L2
normalization. It was chosen in a bake-off against bge-m3, LaBSE and
paraphrase-multilingual-MiniLM for speed, size and an exact ONNX export
(max abs diff 7·10⁻⁸ versus PyTorch), accepting slightly weaker
translation alignment than LaBSE (0.67 vs 0.62 L1 agreement at 256 cells
in the bake-off). §5.3 revisits this choice.

\hypertarget{canonical-space-meaning-lda}{%
\subsubsection{3.2 Canonical space
(meaning-LDA)}\label{canonical-space-meaning-lda}}

With unit vectors \texttt{x}, \texttt{y\ =\ (x\ −\ μ)\ W}: PCA to 384
directions, then the generalized eigenvectors of
\texttt{(S\_total,\ S\_within\ +\ reg)}, where \texttt{S\_within} is the
scatter of translations of the same sentence (FLORES-200 dev, 800
meanings × 12 languages; 7.2k Tatoeba pairs), shrunk toward isotropy
(\texttt{reg\ =\ 0.05}). The top 256 directions are scaled so
within-meaning noise has unit variance. As an ablation, PCA alone gives
34.1\% cross-lingual L1 agreement versus 46.1\% with meaning-LDA (4096
cells).

\hypertarget{quantizer-and-code}{%
\subsubsection{3.3 Quantizer and code}\label{quantizer-and-code}}

Residual k-means with 4 levels (4096·256·256·256) on 77.6k sentences
(Wikipedia sample in 12 languages + FLORES dev). Codes are Crockford
base32 (widths 4/2/2/2) after a published affine permutation, so 85\% of
single-symbol typos are rejected without a checksum.

\hypertarget{confidence}{%
\subsubsection{3.4 Confidence}\label{confidence}}

At each level, a softmax over the 16 nearest centroids' squared
distances, with a temperature fitted by maximum likelihood to predict
the \emph{held-out translation's} symbol; the product over levels is the
prefix confidence. Confidence drives truncation, alternates and the
gate's ``confident portion''.

\hypertarget{adapters}{%
\subsubsection{3.5 Adapters}\label{adapters}}

Other models enter the canonical space through ridge maps fitted on 40k
paired texts (registry: all-MiniLM-L6-v2, bge-small-en-v1.5, gte-small,
paraphrase-multilingual-MiniLM-L12-v2, LaBSE). An unsupervised
vec2vec-style adapter failed at our CPU budget (top-1 retrieval 0\%
versus 96.5\% for ridge on the same texts), so supervised adapters are
the default.

\hypertarget{the-decision-gate}{%
\subsection{4. The decision gate}\label{the-decision-gate}}

Owner review of the first release concluded that in-sample numbers did
not support a ``GPS for meaning'' claim. We therefore fixed the
following gate \textbf{before} any v2 run (D-017). All criteria use
held-out FLORES-200 devtest sentences in 12 languages, never used to fit
the space, codebooks, calibration, labels or adapters.

\begin{longtable}[]{@{}
  >{\raggedright\arraybackslash}p{(\columnwidth - 6\tabcolsep) * \real{0.2500}}
  >{\raggedright\arraybackslash}p{(\columnwidth - 6\tabcolsep) * \real{0.2500}}
  >{\raggedright\arraybackslash}p{(\columnwidth - 6\tabcolsep) * \real{0.2500}}
  >{\raggedright\arraybackslash}p{(\columnwidth - 6\tabcolsep) * \real{0.2500}}@{}}
\toprule\noalign{}
\begin{minipage}[b]{\linewidth}\raggedright
\#
\end{minipage} & \begin{minipage}[b]{\linewidth}\raggedright
Criterion
\end{minipage} & \begin{minipage}[b]{\linewidth}\raggedright
Definition
\end{minipage} & \begin{minipage}[b]{\linewidth}\raggedright
Bar
\end{minipage} \\
\midrule\noalign{}
\endhead
\bottomrule\noalign{}
\endlastfoot
1 & Coverage & share of held-out sentences (all 12 languages) with L1
confidence ≥ 0.7 & ≥ 60\% \\
2 & Translation agreement & among en↔x pairs whose English code has L1
confidence ≥ 0.7, share with the same L1 code (mean over 11 languages) &
≥ 80\% \\
3 & Cross-model agreement & same English text, certified adapter vs
anchor, same L1 code (mean over adapters) & ≥ 60\% \\
4 & Language collapse & per language, share of its sentences in L1 cells
whose training members are ≥ 80\% that language (English excluded) & ≤
5\% for every language \\
\end{longtable}

\textbf{Amendment (D-018).} A tiny-scale dry run ``passed'' all four
criteria with 16 top-level cells and 6 languages. A gate that can be
passed by making the partition coarse says nothing about coordinates. We
amended it: the gate must be met with \textbf{≥ 256 top-level cells}, on
all 12 languages, with the cell count declared before the run. Results
at 32--64 cells are reported as evidence for a \emph{coarse topic}
product only.

\hypertarget{results}{%
\subsection{5. Results}\label{results}}

\hypertarget{full-gate-results}{%
\subsubsection{5.1 Full gate results}\label{full-gate-results}}

\begin{longtable}[]{@{}
  >{\raggedright\arraybackslash}p{(\columnwidth - 12\tabcolsep) * \real{0.1429}}
  >{\raggedright\arraybackslash}p{(\columnwidth - 12\tabcolsep) * \real{0.1429}}
  >{\raggedright\arraybackslash}p{(\columnwidth - 12\tabcolsep) * \real{0.1429}}
  >{\raggedright\arraybackslash}p{(\columnwidth - 12\tabcolsep) * \real{0.1429}}
  >{\raggedright\arraybackslash}p{(\columnwidth - 12\tabcolsep) * \real{0.1429}}
  >{\raggedright\arraybackslash}p{(\columnwidth - 12\tabcolsep) * \real{0.1429}}
  >{\raggedright\arraybackslash}p{(\columnwidth - 12\tabcolsep) * \real{0.1429}}@{}}
\toprule\noalign{}
\begin{minipage}[b]{\linewidth}\raggedright
Build
\end{minipage} & \begin{minipage}[b]{\linewidth}\raggedright
L1 cells
\end{minipage} & \begin{minipage}[b]{\linewidth}\raggedright
Coverage (≥ 60\%)
\end{minipage} & \begin{minipage}[b]{\linewidth}\raggedright
Translation, confident (≥ 80\%)
\end{minipage} & \begin{minipage}[b]{\linewidth}\raggedright
Cross-model (≥ 60\%)
\end{minipage} & \begin{minipage}[b]{\linewidth}\raggedright
Worst collapse (≤ 5\%)
\end{minipage} & \begin{minipage}[b]{\linewidth}\raggedright
Pass
\end{minipage} \\
\midrule\noalign{}
\endhead
\bottomrule\noalign{}
\endlastfoot
v1 frozen release & 4096 & 29.4\% & 65.1\% & 15.7\% & 33.8\% (pt) &
0/4 \\
v1 re-clustered & 1024 & 42.6\% & 69.7\% & 23.1\% & 44.2\% (pt) & 0/4 \\
v1 re-clustered & 256 & 49.2\% & 74.4\% & 32.8\% & 54.6\% (pt) & 0/4 \\
v1 re-clustered & 128 & 54.1\% & 73.6\% & 44.6\% & 0.1\% & 1/4 \\
v1 re-clustered & 64 & 65.0\% & 78.3\% & 57.1\% & 0.0\% & 2/4 \\
v1 re-clustered & 32 & 69.0\% & 79.3\% & 64.5\% & 0.0\% & 3/4 \\
v2-mini (diversified sample) & 4096 & 30.9\% & 65.0\% & 16.3\% & 26.8\%
(pt) & 0/4 \\
v2-mini & 32 & 63.9\% & 78.8\% & 62.3\% & 0.1\% & 3/4 \\
e5-small anchor (20k corpus) & 32 & 55.8\% & 84.2\% & 60.3\%* & 0.0\% &
3/4 \\
LaBSE anchor (20k corpus) & 32 & 66.5\% & 88.7\% & 53.0\%* & 0.0\% &
3/4 \\
LaBSE anchor (20k corpus) & 64 & 64.4\% & 85.3\% & 39.1\%* & 0.0\% &
3/4 \\
\end{longtable}

* Reduced cross-model panel (the other anchor and all-MiniLM-L6-v2).
Sources: \texttt{bench/results/l1-sweep.json},
\texttt{v2-mini-gate.json}, \texttt{v2-mini-4096-gate.json},
\texttt{anchor-check.json}.

\textbf{No configuration passes all four criteria at any granularity,
and none passes more than one criterion at the amended bar of ≥ 256
cells: every build at 256, 1024 or 4096 cells fails all four.} The
closest configurations (3 of 4) all use 32--64 top-level cells, which by
the amendment cannot support a coordinate-system claim.

\hypertarget{the-granularity-trade-off}{%
\subsubsection{5.2 The granularity
trade-off}\label{the-granularity-trade-off}}

Stability is bought with coarseness. Along the v1 re-clustering sweep,
each 4× increase in top-level cells costs roughly 10--25 points on the
gate metrics:

\begin{longtable}[]{@{}llll@{}}
\toprule\noalign{}
Step & Coverage & Translation (confident) & Cross-model \\
\midrule\noalign{}
\endhead
\bottomrule\noalign{}
\endlastfoot
32 → 128 & −14.9 & −5.7 & −20.0 \\
64 → 256 & −15.8 & −3.9 & −24.3 \\
256 → 1024 & −6.6 & −4.7 & −9.8 \\
1024 → 4096 & −13.2 & −4.6 & −7.3 \\
\end{longtable}

Cross-model agreement degrades fastest, then coverage; confident
translation agreement degrades more slowly because the confidence filter
removes the hardest boundary cases (at the price of coverage). The same
pattern holds for LaBSE (32 → 64 cells: −2.1 coverage, −3.4 translation,
−13.9 cross-model). We see no sign of a plateau: finer partitions simply
place more boundaries through regions where equivalent texts land.

\hypertarget{the-anchor-effect-labse-vs-e5-small}{%
\subsubsection{5.3 The anchor effect: LaBSE vs
e5-small}\label{the-anchor-effect-labse-vs-e5-small}}

On the same 20k corpus at 32 cells, LaBSE beats e5-small on coverage
(66.5\% vs 55.8\%, +10.7) and confident translation agreement (88.7\% vs
84.2\%, +4.5), and is the only build that clears the translation bar. It
is worse on cross-model agreement with the reduced panel (53.0\% vs
60.3\%), plausibly because the panel's other models are closer to e5's
geometry. The anchor is therefore a first-order factor, comparable in
size to one 4× granularity step, but it does not change the conclusion:
a better anchor buys roughly one step of resolution, not the six steps
from 32 to 4096 cells. We tested bge-m3 and LaBSE as anchors on a GPU in
§5.5.

\hypertarget{language-collapse}{%
\subsubsection{5.4 Language collapse}\label{language-collapse}}

The first release gave confident top-level codes to only 4.5\% of
Portuguese held-out sentences versus 21--36\% for other languages. Data,
language tags, alignment (en→pt nearest neighbour 99.2\%) and encoding
were correct. The cause: \textbf{52.8\% of Portuguese test sentences
fell into Portuguese-dominated cells} (80--96\% Portuguese training
members; 5 cells hold 53\% of pt test sentences), regardless of topic;
for every other language the share was ≤ 0.2\%
(\texttt{bench/results/language-collapse-small.json}). The Portuguese
Wikipedia sample is unusually homogeneous (mean pairwise cosine 0.84 vs
\textasciitilde0.78, dominated by sports and list stubs), so k-means
spent cells on \emph{language identity} instead of meaning.

Two findings generalize:

\begin{itemize}
\tightlist
\item
  \textbf{Collapse is a function of granularity.} It is absent at ≤ 128
  cells (≤ 0.1\%) and appears abruptly at 256 cells (54.6\% for pt).
  Once the partition is fine enough, residual language signal in the
  anchor becomes a cheaper way for k-means to reduce distortion than
  topic.
\item
  \textbf{The sample explains only part of it.} A diversified
  per-language sample (≤ 40 sentences per title category,
  near-duplicates removed, ≤ 8000 per language) cut the worst collapse
  at 4096 cells only from 33.8\% to 26.8\%. Part of the cause is the
  anchor itself: e5-small keeps a separable Portuguese direction that
  meaning-LDA does not fully remove.
\end{itemize}

Language collapse is invisible in average cross-lingual agreement and
should be reported as a separate metric by any multilingual quantization
work.

\hypertarget{gpu-run-stronger-anchors-at-medium-scale}{%
\subsubsection{5.5 GPU run: stronger anchors at medium
scale}\label{gpu-run-stronger-anchors-at-medium-scale}}

One bounded run (Kaggle T4, 12 languages, \textasciitilde450k sentences,
diversified per-language sample), evaluated on held-out FLORES-200
devtest against the pre-registered bar:

\begin{longtable}[]{@{}
  >{\raggedright\arraybackslash}p{(\columnwidth - 12\tabcolsep) * \real{0.1429}}
  >{\raggedright\arraybackslash}p{(\columnwidth - 12\tabcolsep) * \real{0.1429}}
  >{\raggedright\arraybackslash}p{(\columnwidth - 12\tabcolsep) * \real{0.1429}}
  >{\raggedright\arraybackslash}p{(\columnwidth - 12\tabcolsep) * \real{0.1429}}
  >{\raggedright\arraybackslash}p{(\columnwidth - 12\tabcolsep) * \real{0.1429}}
  >{\raggedright\arraybackslash}p{(\columnwidth - 12\tabcolsep) * \real{0.1429}}
  >{\raggedright\arraybackslash}p{(\columnwidth - 12\tabcolsep) * \real{0.1429}}@{}}
\toprule\noalign{}
\begin{minipage}[b]{\linewidth}\raggedright
anchor
\end{minipage} & \begin{minipage}[b]{\linewidth}\raggedright
L1 regions
\end{minipage} & \begin{minipage}[b]{\linewidth}\raggedright
coverage ≥ 60\%
\end{minipage} & \begin{minipage}[b]{\linewidth}\raggedright
translation (confident) ≥ 80\%
\end{minipage} & \begin{minipage}[b]{\linewidth}\raggedright
cross-model ≥ 60\%
\end{minipage} & \begin{minipage}[b]{\linewidth}\raggedright
collapse ≤ 5\%
\end{minipage} & \begin{minipage}[b]{\linewidth}\raggedright
passes
\end{minipage} \\
\midrule\noalign{}
\endhead
\bottomrule\noalign{}
\endlastfoot
multilingual-e5-small & 256 & 46.3\% FAIL & 74.4\% FAIL & 31.3\% FAIL &
40.9\% FAIL & 0/4 \\
multilingual-e5-small & 1024 & 39.3\% FAIL & 71.6\% FAIL & 24.7\% FAIL &
31.6\% FAIL & 0/4 \\
multilingual-e5-small & 4096 & 26.7\% FAIL & 66.0\% FAIL & 16.2\% FAIL &
34.2\% FAIL & 0/4 \\
LaBSE & 256 & 69.4\% PASS & 83.5\% PASS & 40.1\% FAIL & 0.2\% PASS &
3/4 \\
LaBSE & 1024 & 55.2\% FAIL & 83.9\% PASS & 29.3\% FAIL & 0.1\% PASS &
2/4 \\
LaBSE & 4096 & 46.9\% FAIL & 78.4\% FAIL & 17.7\% FAIL & 0.4\% PASS &
1/4 \\
\textbf{bge-m3} & \textbf{256} & 67.2\% PASS & 88.0\% PASS &
\textbf{50.2\% FAIL} & 0.1\% PASS & 3/4 \\
bge-m3 & 1024 & 62.8\% PASS & 82.6\% PASS & 38.0\% FAIL & 0.1\% PASS &
3/4 \\
bge-m3 & 4096 & 54.5\% FAIL & 78.6\% FAIL & 25.3\% FAIL & 0.2\% PASS &
1/4 \\
\end{longtable}

The stronger anchors remove language collapse (≤ 0.4\% vs 34--41\% for
e5-small) and lift confident translation agreement to 88.0\% (bge-m3,
256 cells). Cross-model agreement does not follow: 50.2\% at best,
falling 50 → 38 → 25\% with each 4× refinement. \textbf{The gate fails},
as we predicted. The bottleneck is not the anchor or the data volume but
the linear adapters between models.

\hypertarget{why-fine-grained-universal-codes-fail}{%
\subsection{6. Why fine-grained universal codes
fail}\label{why-fine-grained-universal-codes-fail}}

Three effects compound. (1) \textbf{Boundary mass:} equivalent texts are
not identical points but small clouds (translation noise is our unit of
distance); the fraction of clouds cut by some cell boundary rises with
the number of cells. (2) \textbf{Adapter error:} a linear adapter adds
its own displacement, which at fine granularity exceeds the cell radius;
hence cross-model agreement falls fastest. (3) \textbf{Residual language
signal:} at fine granularity, k-means can profitably spend cells on
language rather than meaning. Confidence and alternates make (1) visible
but do not remove it; they trade coverage for precision. A code can be
calibrated and honest about its uncertainty and still not be
\emph{universal}.

\hypertarget{what-works}{%
\subsection{7. What works}\label{what-works}}

\hypertarget{coarse-regions}{%
\subsubsection{7.1 Coarse regions}\label{coarse-regions}}

At 32--64 top-level cells, confident translations share a region
78--89\% of the time and there is no language collapse; with the v1
space or the LaBSE anchor, coverage is 64--69\%, and cross-model
agreement is 57--65\% for the v1 space (lower for LaBSE at 64 cells:
39\%, reduced panel). That is a useful frozen, multilingual, shared
topic partition: readable region codes for analytics dashboards, SQL
\texttt{GROUP\ BY}, routing queries to specialised handlers, and
semantic caching keys, with finer levels kept as optional detail
carrying explicit confidence.

\hypertarget{migration-through-stored-canonical-vectors}{%
\subsubsection{7.2 Migration through stored canonical
vectors}\label{migration-through-stored-canonical-vectors}}

The continuous canonical space is a much better hub than the discrete
code. Documents indexed with model X through its adapter and queried
with the anchor reach recall@10 of 96.5\% (all-MiniLM-L6-v2), 99.1\%
(bge-small-en-v1.5), 99.0\% (gte-small) and 98.5\%
(paraphrase-multilingual-MiniLM), versus 100\% for re-embedding, and
98.3\% for LaBSE in the paraphrase setting; using codes alone gives only
10--28\%. Storing the 256-d canonical vector next to each item lets an
organisation change embedding models without re-embedding its corpus,
with N adapters instead of N² pairwise maps.

\hypertarget{other-measured-properties-of-the-frozen-release}{%
\subsubsection{7.3 Other measured properties of the frozen
release}\label{other-measured-properties-of-the-frozen-release}}

Cross-lingual L1 agreement 46.1\% (stable mode 46.7\%), paraphrase L1
agreement 71.6\%; when both codes have L1 confidence ≥ 0.7 (16\% of
pairs), translations agree 88.7\%, at ≥ 0.9 (5\%), 96.5\%. 92\% of L1
cells are used on the training corpus (perplexity 2182 of 4096). Codes
alone are a coarse index: translation retrieval top-1 30\% (raw vectors
99.6\%), SIB-200 topic P@10 16.6\% vs 19.5\%. Adding cross-model scatter
to \texttt{S\_within} helped at 1024 cells but not at 4096, while
lowering cross-lingual agreement; we kept it off.

\hypertarget{limitations}{%
\subsection{8. Limitations}\label{limitations}}

Most experiments are CPU, \texttt{small} scale (≤ 80k fitting sentences,
12 languages, FLORES-200 devtest for evaluation); one GPU run reached
\texttt{medium} scale (\textasciitilde450k sentences, §5.5) and
confirmed the failure at ≥ 256 cells. We did not try nonlinear or
jointly trained adapters. The cross-model panel for the anchor
comparison is reduced. Labels for cells are machine-generated and
filtered, not natively reviewed. FLORES sentences are short, formal and
topically narrow; real traffic is likely harder.

\hypertarget{conclusion}{%
\subsection{9. Conclusion}\label{conclusion}}

Can meaning have universal coordinates? At fine resolution, with today's
small multilingual encoders and linear machinery, our pre-registered
answer is \textbf{no}: no build passes a fixed, held-out gate at ≥ 256
cells, every 4× refinement costs 10--25 points, and fine partitions
start encoding language instead of meaning. At coarse resolution the
answer is a qualified \textbf{yes}: a few dozen shared regions are
stable across languages and reasonably stable across models. And the
continuous space behind the code is a practical hub for moving between
models. We recommend that work on semantic IDs and shared meaning spaces
report held-out, cross-language and cross-model agreement \emph{as a
function of granularity}, together with a language-collapse metric, and
fix the bar before running.

\hypertarget{reproducibility}{%
\subsection{Reproducibility}\label{reproducibility}}

Gate definitions: \texttt{DECISIONS.md} D-017/D-018. Results:
\texttt{bench/results/}. Experimental v2 build:
\texttt{experiments/v2-mini/}. GPU run (§5.5):
\texttt{notebooks/kaggle\_gpu\_gate.ipynb}, or
\texttt{bash\ scripts/gpu\_build\_v2.sh\ -\/-scale\ medium\ -\/-sizes\ "256\ 1024\ 4096"};
results \texttt{bench/results/v2-gate-medium-*.json}. Code, artifacts
and data recipes are Apache-2.0.

\hypertarget{references}{%
\subsection{References}\label{references}}

Adelani et al.~SIB-200, EACL 2024 · Chen et al.~BGE M3, 2024 · Duquenne
et al.~SONAR, 2023 · Feng et al.~LaBSE, ACL 2022 · Grootendorst.
BERTopic, 2022 · Huh et al.~The Platonic Representation Hypothesis, ICML
2024 · Jha et al.~Harnessing the Universal Geometry of Embeddings,
arXiv:2505.12540, 2025 · LCM team. Large Concept Models, 2024 · Lee et
al.~RQ-VAE, CVPR 2022 · NLLB Team. FLORES-200, 2022 · Rajput et
al.~TIGER, NeurIPS 2023 · Singh et al. Semantic IDs for ranking, 2024 ·
Tay et al.~DSI, NeurIPS 2022 · Wang et al.~Multilingual E5, 2024 ·
Zeghidour et al.~SoundStream, 2021 · Bhartṛhari, Vākyapadīya.

\end{document}
